\documentclass[11pt]{article}

\usepackage[a4paper,margin=30mm]{geometry}
\usepackage{amsmath,amssymb,amsthm}
\usepackage[T1]{fontenc}
\usepackage{microtype}
\usepackage[hidelinks]{hyperref}
\hypersetup{
  pdftitle={The Four-Dimensional Hypercube Inequality for Two-Valued Vertex Weights},
  pdfauthor={Xiaojian Zeng}
}

\newtheorem{theorem}{Theorem}
\newtheorem{proposition}[theorem]{Proposition}
\newcommand{\cube}{\{0,1\}^{4}}

\title{The Four-Dimensional Hypercube Inequality\\
for Two-Valued Vertex Weights}
\author{Xiaojian Zeng\\
\small PBC School of Finance, Tsinghua University, Beijing 100190, China\\
\small \texttt{xiaojianzeng@foxmail.com}}
\date{}

\begin{document}

\maketitle

\begin{abstract}
The hypercube inequality in dimension four is a homogeneous degree-eight
inequality in sixteen nonnegative vertex weights.  Its validity on the full
nonnegative orthant is unresolved.  We prove the restricted inequality when
the sixteen weights take at most two distinct values.  For each unordered
nontrivial bipartition of the vertices, substitution of the two values gives
a bivariate polynomial.  Exact enumeration of all $2^{15}-1=32767$
bipartitions shows that all nine coefficients of every such polynomial are
nonnegative.  The enumeration uses integer simplex determinants and is
provided as a reproducible certificate.  Consequently, any strict
counterexample to the unrestricted inequality must have at least three
distinct coordinate values.
\end{abstract}

\section{Statement and context}

Let $V=\cube$.  For a five-element set
$S=\{v_0,\ldots,v_4\}\subset V$, set
\[
 \operatorname{Vol}(S)
 =\left|\det[v_1-v_0\;v_2-v_0\;v_3-v_0\;v_4-v_0]\right|,
\]
the normalized four-dimensional volume of $\operatorname{conv}S$.  For
$i\in\{1,2,3,4\}$ and $\epsilon\in\{0,1\}$, write
\[
 F_{i,\epsilon}=\{v\in V:v_i=\epsilon\}.
\]
These are the eight facets of the four-cube.  For nonnegative weights
$x=(x_v)_{v\in V}$, define
\begin{align*}
 A(x)&=\sum_{S\in\binom V5}\operatorname{Vol}(S)
       \prod_{v\in S}x_v, &
 T(x)&=\sum_{v\in V}x_v,\\
 B(x)&=\prod_{i=1}^{4}\prod_{\epsilon=0}^{1}
       \left(\sum_{v\in F_{i,\epsilon}}x_v\right), &
 H_4(x)&=B(x)-A(x)T(x)^3.
\end{align*}
The four-dimensional hypercube inequality is $H_4(x)\geq0$ on
$\mathbb R_{\geq0}^{16}$.

Averkov, von Dichter, and Soprunov introduced the general hypercube
formulation, proved the preceding case $d=3$, and explicitly left $d\geq4$
unresolved~\cite{AverkovVonDichterSoprunov2026}.  The $d=2$ identity used in
their inductive argument had appeared in an earlier algebraic-combinatorial
proof of a related mixed-volume inequality for three-dimensional
zonoids~\cite{AverkovSoprunov2026}.  The two-entry reduction for auxiliary
matrix rows in
\cite[Section~4.2]{AverkovVonDichterSoprunov2026} is different from requiring
the final vertex weights $x_v$ to take two values.  The result below concerns
only the latter restriction.  It neither proves nor disproves the
unrestricted $d=4$ inequality.  The comparison made here is deliberately
limited to these directly related sources and is not a claim of priority.

\begin{theorem}\label{thm:two-values}
If the coordinates of $x\in\mathbb R_{\geq0}^{16}$ take at most two distinct
values, then $H_4(x)\geq0$.
\end{theorem}

\section{Bivariate reduction}

Fix $U\subset V$ and put $x_v=a$ for $v\in U$ and $x_v=b$ for
$v\in V\setminus U$, where $a,b\geq0$.  Let $m=|U|$ and define
\[
 c_j(U)=\sum_{\substack{S\in\binom V5\\|S\cap U|=j}}
          \operatorname{Vol}(S),
 \qquad
 q_F(U)=|F\cap U|.
\]
Direct grouping gives
\begin{align}
 A(x)&=\sum_{j=0}^{5}c_j(U)a^j b^{5-j},\label{eq:A}\\
 T(x)^3&=(ma+(16-m)b)^3,\label{eq:T}\\
 B(x)&=\prod_F\bigl(q_F(U)a+(8-q_F(U))b\bigr),\label{eq:B}
\end{align}
where the product in \eqref{eq:B} runs over the eight facets.  Hence
\begin{equation}\label{eq:Hexpansion}
 H_4(x)=\sum_{k=0}^{8}h_k(U)a^k b^{8-k}
\end{equation}
with integer coefficients.  More explicitly, if
\[
 \beta_k(U)=
 \sum_{\substack{\mathcal J\subseteq\mathcal F\\|\mathcal J|=k}}
 \prod_{F\in\mathcal J}q_F(U)
 \prod_{F\in\mathcal F\setminus\mathcal J}(8-q_F(U)),
\]
where $\mathcal F$ is the set of facets, then
\begin{equation}\label{eq:hk}
 h_k(U)=\beta_k(U)-
 \sum_{\substack{0\leq j\leq5\\0\leq k-j\leq3}}
 c_j(U)\binom{3}{k-j}m^{k-j}(16-m)^{3-k+j}.
\end{equation}
Thus coefficientwise nonnegativity of \eqref{eq:Hexpansion} is a sufficient
exact certificate for the desired inequality on this two-parameter family.

\begin{proposition}[Exact finite certificate]\label{prop:certificate}
For every nonempty $U\subset V$ that does not contain $(1,1,1,1)$, all nine
integers $h_0(U),\ldots,h_8(U)$ in \eqref{eq:hk} are nonnegative.
\end{proposition}

\begin{proof}
The accompanying program \texttt{two\_level\_certificate.py} implements
\eqref{eq:hk} using only exact integer arithmetic.  It first enumerates all
$\binom{16}{5}=4368$ five-subsets.  Each normalized volume is evaluated by
the Leibniz determinant formula.  The resulting distribution is
\[
\begin{array}{c|rrrr}
\operatorname{Vol}(S)&0&1&2&3\\ \hline
\#S&1360&2672&320&16.
\end{array}
\]
For every mask $1\leq M<2^{15}$, the program takes the corresponding
nonempty subset of the fifteen vertices other than $(1,1,1,1)$, computes the
six numbers $c_j(U)$ and the eight facet intersections, and forms the nine
coefficients by exact polynomial multiplication.  It tests
\[
 (2^{15}-1)\cdot9=294903
\]
integer coefficients.  No coefficient is negative; $2040$ are zero and the
smallest positive coefficient is $1$.

As a consistency check, the weighted sum of the volume distribution is
$2672+2(320)+3(16)=3360$.  Therefore, at the all-ones vector,
\[
 H_4(1,\ldots,1)=8^8-16^3(3360)=3014656,
\]
and the program verifies that $\sum_{k=0}^{8}h_k(U)$ equals this value for
every mask.  The computation terminates with a nonzero exit status if the
volume distribution, the all-ones invariant, or coefficientwise
nonnegativity fails.  Running
\[
 \texttt{python3 evidence/two\_level\_certificate.py}
\]
from the source directory reproduces the calculation.  This is an exhaustive
integer computation, not a floating-point search.
\end{proof}

\begin{proof}[Proof of Theorem~\ref{thm:two-values}]
Suppose first that exactly two values $a$ and $b$ occur.  Their level sets
form an unordered nontrivial bipartition $\{U,V\setminus U\}$.  Exactly one
member of this pair omits the distinguished vertex $(1,1,1,1)$, so it is one
of the $2^{15}-1$ subsets in Proposition~\ref{prop:certificate}.  Exchanging
the two sides only exchanges $a$ and $b$.  By \eqref{eq:Hexpansion},
Proposition~\ref{prop:certificate}, and $a,b\geq0$, every summand
$h_k(U)a^k b^{8-k}$ is nonnegative.  This also covers $a=0$ or $b=0$.

If all coordinates have the same value, choose any enumerated nontrivial
$U$ and set $a=b$; the same conclusion follows.  Thus all assignments with
at most two coordinate values are covered.
\end{proof}

\section{Scope and reproducibility}

The proof is coefficientwise only after restriction to each two-level cone;
it makes no assertion that the full sixteen-variable polynomial has
nonnegative coefficients.  In particular, vectors with three or more
distinct weights are untreated.  The immediate structural consequence of
Theorem~\ref{thm:two-values} is therefore only that a strict counterexample,
if one exists, must use at least three distinct coordinate values.

The certificate is a self-contained Python standard-library program.  During
preparation, its output was also reproduced by a separately written exact
implementation using a different vertex ordering, fraction-free determinant
elimination, and direct expansion of the $2^8$ choices in the facet product.
This cross-check is additional implementation evidence; the mathematical
reduction and the supplied exhaustive certificate are the proof artifacts.

A separate finite-certificate study of ternary Borsuk and kissing-code models
likewise converts geometric compatibility to an integer predicate and checks a
finite-library theorem by exhaustive computation with an independently
reconstructed implementation~\cite[Section~4, Lemma~3 and
Proposition~1]{202608.1658}.  That work concerns different finite models and
provides only a methodological analogue; it does not address the Hypercube
inequality.

An AI research and coding assistant was used to organize the finite
reduction, assist with exact-code implementation and cross-checking, conduct
a bounded literature search, and help draft the manuscript.  Numerical
optimization and language-model output are not treated as proof.  The author
checked the theorem-critical reduction against the exact computation and
takes responsibility for the statement and exposition.

\bibliographystyle{alpha}
\bibliography{references}

\end{document}
